\documentclass{article}
\usepackage[utf8]{inputenc}
\usepackage[spanish]{babel}
\usepackage{tikz}
\usepackage[normalem]{ulem}
\usepackage[left=1cm, right=2cm, top=2cm, bottom=2cm]{geometry}
\usepackage[framemethod=TikZ]{mdframed}
\usetikzlibrary{shapes.geometric, arrows.meta, positioning, fit}

\newmdenv[
  linecolor=black,
  outerlinewidth=1pt,
  roundcorner=5pt,
  innertopmargin=10pt,
  innerbottommargin=10pt,
  backgroundcolor=gray!5
]{actividadbox}

\tikzset{
    blackbox/.style={
        rectangle, 
        draw, 
        fill=white, 
        minimum width=3.5cm, 
        minimum height=2cm, 
        ultra thick, 
        align=center, % Mantenemos el centro
        inner sep=5pt  % Añade un margen interno para que el texto no toque los bordes
    },
    arrow/.style={-Stealth, thick}
}

\begin{document}

\begin{actividadbox}
    \textbf{Actividad:} \sout{Determine} Distinga límites, entorno, entradas y salidas de los siguientes sistemas:
    \begin{enumerate}
        \item La sala de clases IBC 6-29
        \item El televisor de su casa
        \item Una gallina
        \item La cárcel de Valparaíso
        \item La PUCV
        \item Esta clase
    \end{enumerate}
\end{actividadbox}

\section*{Resolución:}

\subsection*{a. Sistema: Sala de clases IBC 6-29}

Podemos distinguir un sistema de Sala de Clases IBC 6-29 como el lugar físico que alberga la interacción entre estudiantes y profesores, donde se imparten conocimientos y se generan aprendizajes.
\vspace{1cm}

\begin{tikzpicture}[node distance=2cm]
    % Caja Negra
    \node[blackbox] (system) {SALA IBC 6-29 \\ \small(Límite: Paredes y Puerta)};
    
    % Entorno
    \node[draw, dashed, inner sep=1cm, fit=(system), label=above left:Entorno: Salas de clases adyacentes] (env) {};

    % Entradas (Inputs) - Corregidas para que apunten HACIA la caja
    \draw[arrow] ([xshift=-2cm, yshift=0.5cm]system.west) node[left, align=right] {Estudiantes/Profesores} -- ([yshift=0.5cm]system.west);
    \draw[arrow] ([xshift=-2cm, yshift=-0.5cm]system.west) node[left] {Información/Guías} -- ([yshift=-0.5cm]system.west);

    % Salidas (Outputs)
    \draw[arrow] (system.east) -- ++(2,0) node[right, align=left] {Estudiantes/Profesores};
\end{tikzpicture}

\subsection*{b. Sistema: Televisor de su casa}

Podemos distinguir un sistema de Televisor como un dispositivo electrónico 
que recibe señales de televisión y energía para producir imágenes y sonidos, 
permitiendo la visualización de contenido multimedia.
\vspace{1cm}

\begin{tikzpicture}[node distance=2cm]
    % Caja Negra
    \node[blackbox] (televisor) {TELEVISOR \\ \small(Límite: Carcasa)};
    
    % Entorno
    \node[draw, dashed, inner sep=1cm, fit=(televisor), label=above left:Entorno: Sala de estar] (env) {};

    % Entradas (Inputs) - Corregidas para que apunten HACIA la caja
    \draw[arrow] ([xshift=-2cm, yshift=0.5cm]televisor.west) node[left, align=right] {Señal de TV / Energía} -- ([yshift=0.5cm]televisor.west);
    \draw[arrow] ([xshift=-2cm, yshift=-0.5cm]televisor.west) node[left] {Señal de Control Remoto} -- ([yshift=-0.5cm]televisor.west);

    % Salidas (Outputs)
    \draw[arrow] (televisor.east) -- ++(2,0) node[right, align=left] {Imagen / Sonido};
\end{tikzpicture}

\subsection*{c. Sistema: Una gallina}

Podemos distinguir un sistema de una gallina como un organismo vivo que 
interactúa con su entorno para sobrevivir, crecer y reproducirse, 
produciendo huevos y otros subproductos.
\vspace{1cm}

\begin{tikzpicture}[node distance=2cm]
    % Caja Negra
    \node[blackbox] (gallina) {GALLINA \\ \small(Límite: Piel/Plumas)};
    
    % Entorno
    \node[draw, dashed, inner sep=1cm, fit=(gallina), label=above left:Entorno: Galpón / Campo] (env) {};

    % Entradas (Inputs) - Corregidas para que apunten HACIA la caja
    \draw[arrow] ([xshift=-2cm, yshift=0.5cm]gallina.west) node[left, align=right] {Maíz / Agua} -- ([yshift=0.5cm]gallina.west);
    \draw[arrow] ([xshift=-2cm, yshift=-0.5cm]gallina.west) node[left] {Oxígeno / Luz Solar} -- ([yshift=-0.5cm]gallina.west);

    % Salidas (Outputs)
    \draw[arrow] (gallina.east) -- ++(2,0) node[right, align=left] {Huevos / Subproductos};
\end{tikzpicture}

\subsection*{d. Sistema: La cárcel de Valparaíso}

Podemos distinguir un sistema de la cárcel de Valparaíso como un lugar físico 
que alberga la interacción entre internos y personal penitenciario, 
donde se imparten normas y se generan comportamientos.
\vspace{1cm}

\begin{tikzpicture}[node distance=2cm]
    % Caja Negra
    \node[blackbox] (carcel) {CÁRCEL DE VALPARAÍSO \\ \small(Límite: Muros y Rejas)};
    
    % Entorno
    \node[draw, dashed, inner sep=1cm, fit=(carcel), label=above left:Entorno: Ciudad / Comunidad] (env) {};

    % Entradas (Inputs) - Corregidas para que apunten HACIA la caja
    \draw[arrow] ([xshift=-2cm, yshift=0.5cm]carcel.west) node[left, align=right] {Internos / Personal} -- ([yshift=0.5cm]carcel.west);
    \draw[arrow] ([xshift=-2cm, yshift=-0.5cm]carcel.west) node[left] {Recursos / Información} -- ([yshift=-0.5cm]carcel.west);

    % Salidas (Outputs)
    \draw[arrow] (carcel.east) -- ++(2,0) node[right, align=left] {Comportamientos \\ / Normas};

\end{tikzpicture}

\subsection*{e. Sistema: La PUCV}

Podemos distinguir un sistema de la PUCV como una institución educativa que 
interactúa con su entorno para formar profesionales, generar conocimiento y 
contribuir al desarrollo social.
\vspace{1cm}

\begin{tikzpicture}[node distance=2cm]
    % Caja Negra
    \node[blackbox] (pucv) {PUCV \\ \small(Límite: Campus / Infraestructura)};
    
    % Entorno
    \node[draw, dashed, inner sep=1cm, fit=(pucv), label=above left:Entorno: Comunidad / Sociedad] (env) {};

    % Entradas (Inputs) - Corregidas para que apunten HACIA la caja
    \draw[arrow] ([xshift=-2cm, yshift=0.5cm]pucv.west) node[left, align=right] {Estudiantes / Profesores} -- ([yshift=0.5cm]pucv.west);
    \draw[arrow] ([xshift=-2cm, yshift=-0.5cm]pucv.west) node[left] {Recursos / Información} -- ([yshift=-0.5cm]pucv.west);

    % Salidas (Outputs)
    \draw[arrow] (pucv.east) -- ++(2,0) node[right, align=left] {Graduados\\ / Conocimiento};
\end{tikzpicture}

\subsection*{f. Sistema: Esta Clase}

Podemos distinguir un sistema de esta clase como un grupo de estudiantes y profesores que 
interactúan en un entorno educativo para aprender, enseñar y desarrollar habilidades.
\vspace{1cm}

\begin{tikzpicture}[node distance=2cm]
    % Caja Negra
    \node[blackbox] (clase) {CLASE \\ \small(Límite: Aula / Plataforma Educativa)};
    
    % Entorno
    \node[draw, dashed, inner sep=1cm, fit=(clase), label=above left:Entorno: Comunidad / Sociedad] (env) {};

    % Entradas (Inputs) - Corregidas para que apunten HACIA la caja
    \draw[arrow] ([xshift=-2cm, yshift=0.5cm]clase.west) node[left, align=right] {Estudiantes / Profesores} -- ([yshift=0.5cm]clase.west);
    \draw[arrow] ([xshift=-2cm, yshift=-0.5cm]clase.west) node[left] {Recursos / Información} -- ([yshift=-0.5cm]clase.west);

    % Salidas (Outputs)
    \draw[arrow] (clase.east) -- ++(2,0) node[right, align=left] {Aprendizaje \\/ Desarrollo de Habilidades};
\end{tikzpicture}

\end{document}